\subsubsection{Ordinary algorithms over arithmetic local henselizations}

The preceding finite-field construction does not depend on field coefficients
once finite presented-ring linear algebra is supplied.  We make that input
explicit over the integers, so that the resulting family includes singular
mixed-characteristic local rings.

\begin{lemma}[Proved: exact linear algorithms over presented integer algebras]
\label{C:integer-algebra-linear-algorithms}
Given a finite presentation
\[
 T=\mathbf Z[z_1,\ldots,z_s]/(F_1,\ldots,F_r),
\]
there are ordinary terminating algorithms for equality and arithmetic in
$T$, membership in finitely generated ideals and submodules of finite free
$T$-modules, finite generators of the kernel of every finite $T$-matrix,
and finite generators of $\operatorname{Ann}_T(a)$ for any $a\in T$.
The same conclusions hold after any supplied finite list of localizations.
All outputs are finite polynomial expressions in the given presentation;
no regularity or torsion-freeness hypothesis is required.
\end{lemma}
\begin{proof}
Put $P=\mathbf Z[z_1,\ldots,z_s]$.
The integers are a coherent strict B\'ezout ring with computable greatest
common divisors, B\'ezout coefficients and divisibility.
Every leading-term submodule of $P^v$ is finitely generated, because $P$ is
Noetherian.  Therefore the module Buchberger algorithm of
\cite[Theorem 4.6 and Algorithm 2.5]{C-GLNY2020} terminates on every finite
list of vectors in $P^v$.  Its integer division algorithm gives membership
and expressions in a computed Gr\"obner basis.  The integer Schreyer
algorithm \cite[Theorem 6.5 and Algorithm 2.8]{C-GLNY2020} computes finite
generators of the syzygies of such a basis.  These are the precise source
statements used here; no finite projective resolution or bound on its
length is needed.

We record explicitly how to recover syzygies of an arbitrary input matrix,
since the input columns need not already be a Gr\"obner basis.
Let $F:P^u\to P^v$ be that matrix and $G:P^a\to P^v$ a computed
Gr\"obner basis for its image.
Tracking the elementary operations and performing membership division gives
matrices $A:P^a\to P^u$ and $B:P^u\to P^a$ satisfying
\[
                     G=FA,\qquad F=GB.                 \tag{C.3.IH1}
\]
If $w_1,\ldots,w_b$ generate $\ker G$, then the columns of $1-AB$
together with $Aw_1,\ldots,Aw_b$ generate $\ker F$.
They lie in that kernel by (C.3.IH1).  Conversely, for $Fx=0$ we have
$GBx=0$ and
$x=(1-AB)x+A(Bx)$, which proves generation.
Zero matrices and zero columns are handled directly or retained as their
obvious basis syzygies.

For a matrix $M:T^u\to T^v$, choose a polynomial lift $\widetilde M$.
Apply the preceding procedure to the polynomial matrix with columns
\[
 \bigl[\,\widetilde M\ \big|\ -F_j e_l\ (1\leq j\leq r,
                                     1\leq l\leq v)\,\bigr].       \tag{C.3.IH2}
\]
Project its syzygy generators onto their first $u$ coordinates and reduce
modulo $(F_1,\ldots,F_r)$.
They generate the full $T$-kernel: a lift $x\in P^u$ represents an element
of $\ker M$ exactly when $\widetilde Mx$ is a combination of the
$F_j e_l$, and such a combination completes $x$ to a syzygy of (C.3.IH2).
The same observation reduces submodule membership over $T$ to membership
in the $P$-submodule generated by lifted columns and all $F_j e_l$.
Equality is the special case of membership in the presentation ideal,
and annihilators are kernels of the one-by-one multiplication matrices.
Finally, inverting $g\in T$ adds a variable $y$ and the relation $yg-1$;
finitely many localizations therefore remain in the stated class.
Every step uses finite integer and polynomial data, and all the invoked
algorithms terminate.
\end{proof}

\begin{theorem}[Proved: every arithmetic local henselization with finite odd residue]
\label{C:arithmetic-henselization-WG}
Let $S=\mathbf Z[z_1,\ldots,z_s]/I$ be a supplied finite presentation,
and let $S\twoheadrightarrow\kappa$ be a supplied surjection onto a finite
field of odd characteristic, specified by its finite arithmetic and the
images of the generators.  Put $\mathfrak m=\ker(S\to\kappa)$ and
$R=(S_{\mathfrak m})^h$.
Then $R$ has finite names with computable exact arithmetic, equality,
residues and unit inversion.  An ordinary terminating algorithm computes
finite generators of the full kernel of every finite matrix over $R$.

For every $m\geq0$ these data give ordinary terminating algorithms for
the complete classification and all the operations of
Theorem~\ref{C:effective-henselian-DVR-WG}: literal canonical
indecomposable rational $O_m/R$ indices, their duality and hermitian types,
the explicit direct-sum description (C.3.HD2) of the split
$\mathrm{WG}(BO_m/R)$, actual isometry and group equality, tensor products,
and exterior operations.  No arithmetic, polynomial factorization,
approximation, branch-selection or stabilization oracle is part of the
input.  No regularity, reducedness, domain, equicharacteristic, DVR or
positive-dimensional assumption is imposed.
The canonical coefficient-name set is a decidable subset of $\mathbf N$;
this formulation includes finite $R$.
\end{theorem}
\begin{proof}
We give the passage from presented integer algebras to the actual
henselization, and then identify every input needed by the representation
and form algorithms.

\emph{Names and common pointed neighborhoods.}
Use standard pointed \'etale neighborhoods over $S$.
Concretely choose $b\in S$ with $\bar b\ne0$ in $\kappa$, a monic
$f\in S_b[U]$, a root $a\in\kappa$ of $\bar f$ with
$\bar f'(a)\ne0$, and $h\in S_b[U]$ with $\bar h(a)\ne0$.
Set
\[
 T=S_b[U,V]/\bigl(f(U),\,Vf'(U)h(U)-1\bigr),\qquad
 \rho(U)=a,\quad \rho(V)=(\bar f'(a)\bar h(a))^{-1}.       \tag{C.3.IH3}
\]
The map $\rho:T\twoheadrightarrow\kappa$ extends the given one on $S$.
All validity tests in this finite presentation are decidable, and $T$ is
a finitely presented \'etale $S$-algebra.
Include the inverse of $b$ as another variable and relation: $T$ is then
literally a finitely presented integer algebra of the preceding lemma.
A raw name is such a neighborhood together with a polynomial element of $T$.
For intermediate computations allow finite tensor products of these
neighborhoods over $S$, with the product residue map.
They are again pointed finitely presented \'etale algebras and have explicit
finite integer presentations.

Every element of $R$ has a raw name.
Indeed henselization is the filtered colimit of pointed \'etale neighborhoods
with unchanged residue field, and the standard \'etale neighborhood lemma
makes each relevant presentation monic after shrinking at its point
\cite[Tag 0BSK, Lemma 10.155.1; Tag 00UE, Proposition 10.144.4]
 {C-StacksHenselianEffective}.
The finitely many coefficients over $S_{\mathfrak m}$ lie in one $S_b$,
and the remaining shrinking is the factor $h$ in (C.3.IH3).
Any finite set of names maps to their tensor neighborhood, which is
sufficient for all arithmetic and matrix computations below.

For such a neighborhood put $\mathfrak q=\ker\rho$.
There is a canonical pointed identification
\[
                      (T_{\mathfrak q})^h=R.             \tag{C.3.IH4}
\]
The local map $S_{\mathfrak m}\to T_{\mathfrak q}$ is a localization
of an \'etale map with the same residue field.  Its unique pointed lift
to the henselian ring $R$, and the universal property of the two
henselizations, give maps both ways in (C.3.IH4).
One composite is the unique map of the henselization fixing
$S_{\mathfrak m}$; the other fixes $T$ by uniqueness of its pointed
\'etale lift, and hence is the identity after henselization.
Consequently $T\to R$ is flat and $T_{\mathfrak q}\to R$ is faithfully
flat, using the flat local henselization map
\cite[Tag 07QL, Lemma 15.46.1(1); proof of Lemma 15.46.3(c)]
 {C-StacksHenselianEffective}.

\emph{Decidable equality, not formal equality.}
For $z\in T$ we have
\[
 \begin{split}
 z=0\text{ in }R
 &\ \Longleftrightarrow\ z=0\text{ in }T_{\mathfrak q}\\
 &\ \Longleftrightarrow\operatorname{Ann}_T(z)
                                     \not\subseteq\mathfrak q.
 \end{split}                                             \tag{C.3.IH5}
\]
Compute finite generators of this annihilator by
Lemma~\ref{C:integer-algebra-linear-algorithms}.
The last condition is equivalent to one generator having nonzero image
under the supplied finite map $\rho$.
Thus it is decidable even when $T$ has torsion or zero divisors.
For two names pass to their tensor neighborhood and apply (C.3.IH5) to
the difference.  Compute sums and products in the same way.
The residue is given by $\rho$; an element is a unit exactly when its
residue is nonzero, and its inverse is named by adjoining an inverse
variable to $T$, another pointed \'etale localization.

Enumerate all raw finite names by natural-number codes.
Their validity is decidable, including the finite-field tests in
(C.3.IH3).  To replace an internally constructed name by a name in this
fixed enumeration, search for an equal raw name using (C.3.IH5).
Coverage proves termination.
Retain a valid code precisely when no smaller valid code has the same
value.  This is a finite predecessor test, so the retained codes form a
decidable set of unique names.  All searches may run through raw codes;
they never ask for a nonexistent next distinct element when $R$ is finite.
Lifts of the finitely many elements of $\kappa$ can be found by enumerating
$S$ and using its given surjection.

\emph{Full matrix kernels.}
Put the entries of a finite matrix $M$ over $R$ in a common tensor
neighborhood $T$ and choose the resulting matrix $M_T$.
The integer-algebra lemma computes finite generators of $\ker M_T$.
Flatness in (C.3.IH4) proves
\[
                   \ker M=(\ker M_T)\otimes_T R.         \tag{C.3.IH6}
\]
Their images are therefore generators of the actual full $R$-kernel.
No finite-adic image, unproved bound or completion solution is substituted
for that kernel.

\emph{Termination of classification and form operations.}
The ring $R$ is Noetherian henselian local with residue $\kappa$;
Noetherianity is preserved by henselization
\cite[Tag 06LJ, Lemma 15.46.3]{C-StacksHenselianEffective}.
Since the residue characteristic is odd, $2$ is a unit in $R$.
The proof of Theorem~\ref{C:finite-field-henselization-WG} now applies
with its polynomial linear-algebra paragraph replaced by
Lemma~\ref{C:integer-algebra-linear-algorithms}; we spell out why these
are all its inputs.

The integral finite free coefficient spaces
of Lemma~\ref{C:integral-orthogonal-coalgebras} turn coactions,
intertwiners and invariant bilinear forms into finitely many exact
coefficient equations over $R$.
The kernels in (C.3.IH6) give finite generators of every actual Hom and
form module, and their residues have a finite computable image in a
matrix space over $\kappa$.
An intertwiner or Gram matrix is invertible exactly when its determinant
has nonzero residue.  Thus these finite residue images decide whether
an invertible actual solution exists.

For every finite endomorphism $R$-algebra $E$, its image in the finite
residue matrix algebra is computable.  The kernel lies in
$\operatorname{Jac}(E)$, and
Lemma~\ref{C:henselian-finite-algebra-lifting} lifts idempotents through
that kernel.  Thus finite residue calculations decide splittings,
and enumeration of exact endomorphisms finds every certified lift.
Splitting decreases positive ranks.  The image of an idempotent is free
because $R$ is local; independent residue columns lift a basis, with
independence certified by a unit minor and generation by Nakayama.
Least finite code among all isomorphic predecessors gives the literal
indecomposable indices, by the already decidable invertible-Hom test.

The same lifting lemma identifies hermitian isometry with the corresponding
finite semisimple residue calculation.  Finite division algebras there are
finite fields.  Their symmetric, alternating, unitary and dual-pair cases
therefore give exactly (C.3.HD2), with its computed generators and finite
coordinate equality tests.  Every positive search for an idempotent,
reference form or isometry follows a finite existence decision and hence
terminates.  Tensor matrices, exterior minors and the finite virtual
exterior formulas use these same operations.
These arguments prove every stated ordinary algorithm under precisely the
supplied finite presentation and finite residue input.

For example the theorem applies, for every odd prime $p$, to
\[
 \bigl(\mathbf Z[x,y]/(xy-p^2)\bigr)_{(p,x,y)}^h,
 \qquad
 \bigl(\mathbf Z[x]/(px,x^2)\bigr)_{(p,x)}^h.              \tag{C.3.IH7}
\]
The first local ring is singular: its dimension is two and its maximal
ideal has three independent generators modulo its square, properties
preserved by the identified henselian completion.
In the second ring the nonzero nilpotent $x$ survives localization and
henselization: in the presented ring its annihilator is $(p,x)$, so
localization at that ideal cannot kill it, and henselization is faithfully
flat.  Both examples have characteristic zero and finite odd residue.
They illustrate features allowed by the theorem, rather than additional
assumptions used by its proof.
\end{proof}
